\FloatBarrier
\section{Valley States and Interface Coupling}
\label{app:valley-states}
\label{sec:device-orbital-valley}

Interface coupling can mix the valleys retained in a dot and correlate
their spatial envelopes. This appendix develops the multivalley block
Hamiltonian of Eq.~\eqref{eq:dot-coupled-valley-equations}, then projects it
onto one common orbital of two equivalent valleys. This projection
gives the two-level Hamiltonian and valley eigenstates in
Eqs.~\eqref{eq:dot-valley-doublet} and~\eqref{eq:dot-valley-eigenstates},
and states the conditions under which their splitting is a useful
description of the valley gap.

\subsection{Coupled valley envelopes}
\label{sec:device-coupled-valley-envelopes}

The silicon conduction band has several energy minima at different
wave vectors $\mathbf k_\nu$. The Bloch states near each minimum
span a low-energy space $\mathcal H_\nu$. Each space contains many
wave vectors and therefore many possible spatial wave packets.
The retained wave-vector regions are disjoint, so these valley
spaces are mutually orthogonal.

Within the effective-mass description, the expansion coefficients
for wave vectors near one minimum are represented in position space
by a slowly varying envelope $F_\nu(\mathbf r)$. Retaining $N_v$
valleys gives the direct-sum space and its envelope representation:
\begin{equation}
 \mathcal H_{\mathrm{mv}}
       =\bigoplus_{\nu=1}^{N_v}\mathcal H_\nu,
 \qquad
 \mathbf F(\mathbf r)
       =\begin{pmatrix}F_1(\mathbf r)\\ \vdots\\F_{N_v}(\mathbf r)\end{pmatrix}.
 \label{eq:device-valley-space}
\end{equation}
Each component is a function describing the spatial state within
one valley. The state is normalized by
$\sum_\nu\int |F_\nu|^2\,d^3r=1$.

Let $h_{0,\nu}$ denote the single-valley orbital Hamiltonian
$h_\nu$ of Eq.~\eqref{eq:device-envelope-hamiltonian}. It acts
within the $\nu$th component of the direct-sum space.
An abrupt interface can transfer amplitude between
valleys and mix their spatial states~\cite{Friesen2010}.
We describe transfer from valley $\nu'$ to valley $\nu$ by
an operator $\hat v_{\nu\nu'}$. The full Hamiltonian is
\begingroup
\let\valleyblocklabel\label
\begin{equation}
 h_{\mathrm{mv}}=
 \begin{pmatrix}
  h_{0,1} & \hat v_{12} & \cdots & \hat v_{1,N_v}\\
  \hat v_{21} & h_{0,2} & \cdots & \hat v_{2,N_v}\\
  \vdots & \vdots & \ddots & \vdots\\
  \hat v_{N_v,1} & \hat v_{N_v,2} & \cdots & h_{0,N_v}
 \end{pmatrix}.
 \label{eq:device-multivalley-hamiltonian}
 % Use the ordinary label command for a second alias in this amsmath display.
 \begingroup\let\label\valleyblocklabel
 \label{eq:device-coupled-valley-equations}\endgroup
\end{equation}
\endgroup
Hermiticity requires
$\hat v_{\nu'\nu}=\hat v_{\nu\nu'}^\dagger$.
All diagonal terms are included in $h_{0,\nu}$; the coupling
operators $\hat v_{\nu\nu'}$ connect distinct valleys.

We use the local multivalley effective-mass model, in which each
off-diagonal operator multiplies an envelope by a position-dependent
coupling function~\cite{QTCADMultivalley}:
\[
 (\hat v_{\nu\nu'}F_{\nu'})(\mathbf r)
       =v_{\nu\nu'}(\mathbf r)F_{\nu'}(\mathbf r),
 \qquad
 v_{\nu'\nu}(\mathbf r)=v_{\nu\nu'}^*(\mathbf r).
\]
Thus components at the same position are coupled. The functions
$v_{\nu\nu'}$ encode the microscopic interface structure and are
additional inputs to this effective model.

\subsection{A two-valley doublet within one orbital}
\label{sec:device-valley-doublet-projection}

Consider the two lowest, equivalent valleys with orthogonal labels
$\ket{\nu_1}$ and $\ket{\nu_2}$. Suppose that their uncoupled
orbital blocks $h_{0,\nu_1}$ and $h_{0,\nu_2}$ share a normalized
ground orbital $\phi_0$ with energy $\epsilon_0$.
Retain the states $\ket{\phi_0}\otimes\ket{\nu_1}$ and
$\ket{\phi_0}\otimes\ket{\nu_2}$ as the lowest orbital--valley
pair, with all other states higher in energy. The approximation requires
coupling matrix elements to other orbitals and valleys to be small
compared with their energy separations from this pair. It need not
be small compared with the separation within the pair, which is
zero before the valleys are coupled.

Write $v_{12}(\mathbf r)=v_{\nu_1\nu_2}(\mathbf r)$.
Projection of the local interface coupling onto the common orbital gives
\begin{equation}
 v_0=\bra{\phi_0}v_{12}\ket{\phi_0}
     =\int d^3r\,\phi_0^*(\mathbf r)
                    v_{12}(\mathbf r)\phi_0(\mathbf r).
 \label{eq:device-valley-projected-coupling}
\end{equation}
Hermiticity supplies the conjugate matrix element in the reverse
direction. Define
$P_0=\ket{\phi_0}\bra{\phi_0}\otimes
(\ket{\nu_1}\bra{\nu_1}+\ket{\nu_2}\bra{\nu_2})$.
In the ordered valley basis $(\ket{\nu_1},\ket{\nu_2})$, with
the common orbital suppressed, the restriction of
$P_0h_{\mathrm{mv}}P_0$ is
\begin{equation}
 H_{\mathrm v}=\epsilon_0 I+
       \begin{pmatrix}0&v_0\\v_0^*&0\end{pmatrix}.
 \label{eq:device-valley-doublet}
\end{equation}
The characteristic equation is
$(\epsilon-\epsilon_0)^2-|v_0|^2=0$.
For $v_0\ne0$, write $v_0=|v_0|e^{i\theta_{\mathrm v}}$.
The energies, normalized valley eigenstates, and their gap are
\begin{equation}
 \begin{aligned}
 \epsilon_\pm&=\epsilon_0\pm|v_0|,\\
 \ket{v_\pm}&=\frac{\ket{\nu_1}
                 \pm e^{-i\theta_{\mathrm v}}\ket{\nu_2}}{\sqrt2},\\
 \Delta_{\mathrm{val}}&=\epsilon_+-\epsilon_-=2|v_0|.
 \end{aligned}
 \label{eq:device-valley-eigenstates}
\end{equation}
Thus the interface selects two orthogonal valley superpositions
and separates their energies. If $v_0=0$, the pair remains
degenerate and $\theta_{\mathrm v}$ is undefined; any orthonormal
pair within this valley space is then an eigenbasis.

Within this projection, the lowest orbital--valley state is
$\ket{\chi_0}\simeq\ket{\phi_0}\otimes\ket{v_-}$ for $v_0\ne0$.
The common orbital and equivalent diagonal energies are assumptions
of this two-valley limit. If interface coupling appreciably mixes
other orbitals or valleys, the full block Hamiltonian in
Eq.~\eqref{eq:device-coupled-valley-equations} must be used instead.
Its lowest state can have different spatial envelopes in its
valley components, so $\ket{\chi_0}$ need not factor into one
orbital and one valley state.
